% TOP_PROBLEM: {"id": 4700014, "title": "Maximum number of critical periods", "category_id": 11, "category": {"id": 11, "name": "dynamical_systems", "display_name": "Dynamical Systems", "description": "Problems about long-term behavior of deterministic systems, Hamiltonian dynamics, and periodic orbits.", "slug": "dynamical-systems", "order_index": 11, "created_at": "2026-07-31T15:26:25.671Z"}, "status": "open", "problem_number": "AMR-046-0014", "set_id": 13}
% FIRST_POSTED: 2026-10-02
% ENTRY_KIND: statement_only
% UnsolvedMath ID 4700014; source catalogue code AMR-046-0014
% !TeX program = lualatex
% Definition and sources only; this file is not an LLM solution attempt.
\documentclass[11pt,a4paper]{article}
\usepackage[margin=25mm]{geometry}
\usepackage{fontspec}
\setmainfont{lmroman10-regular.otf}[BoldFont=lmroman10-bold.otf,
 ItalicFont=lmroman10-italic.otf,BoldItalicFont=lmroman10-bolditalic.otf]
\usepackage{amsmath,amssymb,mathtools}
\usepackage{unicode-math}
\setmathfont{latinmodern-math.otf}
\AtBeginDocument{\let\setminus\smallsetminus}
\usepackage{xurl,hyperref}
\hypersetup{colorlinks=true,urlcolor=blue}
\setcounter{secnumdepth}{0}
\setlength{\parindent}{0pt}
\setlength{\parskip}{5pt}
\setlength{\emergencystretch}{3em}
\begin{document}
\section*{Maximum number of critical periods}\label{4700014}
{\small UnsolvedMath ID 4700014 · AMR-046-0014 · Dynamical Systems · AMR Open Problem Lists}
\subsection{Definitions and mathematical statement}
A planar polynomial differential system of degree at most $n$ has the form
$\dot x=P(x,y)$, $\dot y=Q(x,y)$ with real polynomials of total degree at
most $n$. For any center, its period annulus is the maximal surrounding
family of nonconstant periodic orbits. Choosing a regular transverse
coordinate $s$ on this annulus gives the least-period function $T(s)$.
An isolated zero of $T'(s)$ is a critical period; the definition is
independent of the chosen regular transverse coordinate.

Let $\mathcal T(n)$ be the supremum of the numbers of such critical periods
on a period annulus, over all degree-at-most-$n$ systems and their centers.
The problem asks whether there are constants $C>0$ and $n_0$ such that
\[
 \mathcal T(n)\ge C n^2\log n\qquad(n\ge n_0).
\]
Equivalently, can one construct, for every sufficiently large degree, a
system and a center whose surrounding annulus has at least a constant
multiple of $n^2\log n$ distinct stationary periods?

The logarithm is natural; changing its base only changes $C$. The same
positive constant must work for all large $n$. Critical periods are counted
as distinct stationary orbits, without multiplicity, and endpoints of the
annulus are excluded. An isochronous annulus contributes no isolated
critical periods. The intended target is an improvement over quadratic
growth by a logarithmic factor, not an upper estimate on $\mathcal T(n)$
or merely a degree-dependent existence result for a single critical period.
\subsection{Short English statement}
Can degree-$n$ polynomial centers have at least a constant times $n^2\log n$ distinct critical periods?
\subsection{Sources}
\begin{itemize}
\item[S1] ULAM.AI and the UnsolvedMath Contributors, supplied problems.json, record 4700014 (AMR-046-0014), AMR Open Problem Lists. Catalogue identity and source transcription; CC BY 4.0.
\url{https://huggingface.co/datasets/ulamai/UnsolvedMath}.
\item[S2] Gasull - Some open problems in low dimensional dynamical systems (2020), Problem environment 14: Maximum number of critical periods. Original source indexed by AMR-046-0014.
\url{https://arxiv.org/abs/2012.02524}.
\end{itemize}
\end{document}
